%! AP Statistics — Unit 2, Lesson 2.5 — Group Activity
\documentclass[10pt]{article}
\usepackage{apstats-article}
\usepackage{apstats-boxes}

\begin{document}

\pageheader{Unit 2, Lesson 2.5}{Group Activity}
\namepartnerperiod

\vspace{0.14in}

\begin{scenariobox}[Context]{sky}
\small
Your group will calculate and interpret the correlation coefficient $r$ for a real-world
dataset. Each tier uses a different dataset --- compare your $r$ values and interpretations
with other groups at the end.
\end{scenariobox}

\vspace{0.1in}

% ── TIER R ────────────────────────────────────────────────────────────────────
\begin{tcolorbox}[colback=white, colframe=black!40, boxrule=0.8pt, arc=2mm,
    left=4mm, right=4mm, top=3mm, bottom=3mm,
    title={\textbf{Tier R --- Remediate}}, fonttitle=\bfseries, halign title=center]
\small

\textbf{Dataset: Hours of practice ($x$) vs.\ free-throw percentage ($y$) for 10 players}

\vspace{6pt}
\renewcommand{\arraystretch}{1.4}
\begin{tabularx}{\linewidth}{@{} l *{10}{C{0.82cm}} @{}}
\rowcolor{sky}
\textbf{Practice (hr)} & 1 & 2 & 3 & 4 & 5 & 6 & 7 & 8 & 9 & 10 \\
\textbf{FT \%}         & 55 & 60 & 63 & 67 & 70 & 73 & 76 & 80 & 83 & 86 \\
\end{tabularx}

\vspace{4pt}
\textbf{Summary statistics (pre-computed):}
$\bar{x} = 5.5$,\; $\bar{y} = 71.3$,\; $s_x = 3.03$,\; $s_y = 10.19$

\vspace{10pt}
\begin{enumerate}[leftmargin=1.6em, itemsep=0.16in]
  \item Use technology to compute $r$. \quad $r = \blank{2cm}$
  \item Write a complete interpretation of $r$ in context.\\[6pt]
        \writeline\\[2pt]\writeline\\[2pt]\writeline
  \item Does this mean that practicing more \emph{causes} a player to make more
        free throws? Explain.\\[6pt]
        \writeline\\[2pt]\writeline
\end{enumerate}
\end{tcolorbox}

\vspace{0.1in}

% ── TIER A ────────────────────────────────────────────────────────────────────
\begin{tcolorbox}[colback=white, colframe=black!40, boxrule=0.8pt, arc=2mm,
    left=4mm, right=4mm, top=3mm, bottom=3mm,
    title={\textbf{Tier A --- Approaching Proficiency}}, fonttitle=\bfseries, halign title=center]
\small

\textbf{Dataset: Number of absences ($x$) vs.\ final exam score ($y$) for 12 students}

\vspace{6pt}
\renewcommand{\arraystretch}{1.4}
\begin{tabularx}{\linewidth}{@{} l *{12}{C{0.72cm}} @{}}
\rowcolor{sky}
\textbf{Absences} & 0 & 1 & 1 & 2 & 3 & 3 & 4 & 5 & 5 & 6 & 7 & 8 \\
\textbf{Exam}     & 96 & 91 & 88 & 84 & 79 & 76 & 70 & 65 & 62 & 58 & 51 & 44 \\
\end{tabularx}

\vspace{10pt}
\begin{enumerate}[leftmargin=1.6em, itemsep=0.16in]
  \item Use technology to compute $r$. \quad $r = \blank{2cm}$
  \item Write a complete interpretation of $r$ in context. Include direction, strength,
        and the word ``linear.''\\[6pt]
        \writeline\\[2pt]\writeline\\[2pt]\writeline
  \item A student says, ``This proves that missing class \emph{causes} you to fail.''
        Evaluate this claim. Suggest one possible lurking variable.\\[6pt]
        \writeline\\[2pt]\writeline\\[2pt]\writeline
  \item Suppose one student had 10 absences but scored 90. How might this outlier affect $r$?
        Explain without recalculating.\\[6pt]
        \writeline\\[2pt]\writeline
\end{enumerate}
\end{tcolorbox}

\vspace{0.1in}

% ── TIER E ────────────────────────────────────────────────────────────────────
\begin{tcolorbox}[colback=white, colframe=black!40, boxrule=0.8pt, arc=2mm,
    left=4mm, right=4mm, top=3mm, bottom=3mm,
    title={\textbf{Tier E --- Extension}}, fonttitle=\bfseries, halign title=center]
\small

\textbf{Dataset: Latitude (°N, $x$) vs.\ average July temperature (°F, $y$) for 10 U.S.\ cities}

\vspace{6pt}
\renewcommand{\arraystretch}{1.4}
\begin{tabularx}{\linewidth}{@{} l *{10}{C{0.82cm}} @{}}
\rowcolor{sky}
\textbf{Latitude}  & 25 & 30 & 33 & 35 & 38 & 40 & 42 & 44 & 47 & 48 \\
\textbf{Temp (°F)} & 90 & 88 & 84 & 82 & 76 & 73 & 71 & 67 & 65 & 63 \\
\end{tabularx}

\vspace{10pt}
\begin{enumerate}[leftmargin=1.6em, itemsep=0.16in]
  \item Use technology to compute $r$. \quad $r = \blank{2cm}$
  \item Write a complete interpretation of $r$ in context.\\[6pt]
        \writeline\\[2pt]\writeline\\[2pt]\writeline
  \item Can you conclude that being at a higher latitude \emph{causes} lower temperatures?
        What is different about this situation compared to a typical lurking-variable example?\\[6pt]
        \writeline\\[2pt]\writeline\\[2pt]\writeline
  \item Now convert the temperatures to Celsius: $C = \tfrac{5}{9}(F - 32)$.
        Without recalculating, predict the new $r$. Explain, then verify with technology.\\[6pt]
        Prediction: \blank{2cm} \quad Reason: \blank{8cm}\\[4pt]
        Technology gives: $r = \blank{2cm}$. Was your prediction correct?
\end{enumerate}
\end{tcolorbox}

\end{document}
